% =============================================================================
% introduction.tex — Paper Introduction (draft v0.3, 2026-09-03)
% Target: Computers & Chemical Engineering (Sampat/Tominac home journal)
% Length target: ~1,300 words. All citations verified 2026-09-03
% (see bib_doi_audit_report.md in 03_Data/WI_Working_Files).
% =============================================================================

\section{Introduction}

Biochar, the carbon-rich solid produced by pyrolyzing biomass under oxygen-limited
conditions, is one of the most scalable engineered carbon dioxide removal (CDR)
pathways available today \cite{woolf2010, lehmann2021biochar}. Applying biochar to
soils simultaneously sequesters photosynthetically fixed carbon on centennial
timescales and delivers agronomic co-benefits, including reduced nitrous oxide
emissions and improved water retention \cite{cayuela2014, smith2016}. Recent
assessments estimate that biochar could remove 0.3--2 Gt CO$_2$ per year globally
by mid-century \cite{woolf2021est}, and the 2023 Billion-Ton Report identifies
agricultural residues and forestry wastes in the United States, more than a
billion dry tonnes per year under mature-market scenarios, as an abundant
feedstock base \cite{bt23}. In parallel, voluntary carbon markets have matured to
the point where biochar-based removal credits trade at \$100--250 per tonne
CO$_2$e under methodologies such as Verra VM0044 and Puro.earth CORC, with
corporate buyers purchasing hundreds of thousands of tonnes annually
\cite{cdrfyi2024, verraVM0044}.

Realizing this potential requires coordinating biomass collection, drying,
pyrolysis, biochar distribution, and soil application across hundreds of
facilities and counties, a classical supply chain network design problem with an
unusually strong coupling between economics and carbon accounting. Unlike
conventional biomass-to-energy systems, a biochar system's dominant value stream
under carbon policy is the removal credit itself, which accrues per tonne of
stable carbon rather than per unit of energy sold. This coupling raises
questions that existing models do not answer: ``How should carbon policies be
formulated inside a supply chain optimization model so that they genuinely
change facility locations, technology choices, and production volumes?'' and
``What is the shadow price of carbon that emerges from such a system, and
how does it compare with observed market prices?''

Three strands of literature frame these questions. First, biochar techno-economic
and life-cycle assessments have established the unit economics of slow pyrolysis
and the permanence accounting of biochar carbon
\cite{roberts2010, woolf2010, nematian2021}; these studies, however, treat a
single plant or a stylized regional average and do not optimize network structure.
Second, biomass and negative-emission supply chain optimization has developed
spatially explicit mixed-integer models for biochar source--sink networks
\cite{tan2016biochar, belmonte2018}, BECCS chains \cite{negri2021}, and
multi-product waste networks \cite{sampat2019}. The market-clearing framework of
Sampat et al.~\cite{sampat2019} and the spatio-temporal economic analysis of
Tominac et al.~\cite{tominac2022} are particularly relevant: they show that the
dual variables of mass-balance constraints in an optimization model are the
``inherent values'' of products at each location, i.e., equilibrium prices
under perfect coordination. Third, the green supply chain literature has compared
carbon taxes, cap-and-trade, and carbon credit policies in generic production and
inventory networks \cite{benjaafar2013, zakeri2015, chaabane2012, elhedhli2012},
demonstrating that a linear tax shifts the optimal solution only when emissions
intensity varies endogenously across decisions; it is otherwise a constant
shift of the objective.

This paper addresses the gap between these strands: no existing
study (i) embeds carbon policy inside a multi-product, price-responsive
supply chain network for biochar in the market-clearing/dual framework, (ii)
derives the shadow carbon price of the system as the dual of an
emission constraint and benchmarks it against real market prices, and (iii)
compares three canonical policy architectures (baseline-and-credit,
cap-and-trade, and a tiered tax with removal crediting and allowance
allocation) on identical network, data, and demand structures. We fill this gap
with a case study of Wisconsin, a state whose 14 million acres of farmland, 46\%
forest cover, and dairy-dominated agriculture make it representative of the
Midwestern biomass landscape.

The contributions of this paper are threefold. First, we formulate a
multi-product biochar supply chain network design model, with 72 county nodes,
13 feedstocks, drying plus slow pyrolysis at two temperatures (300 and
500~$^\circ$C), and a three-segment price-responsive biochar demand curve, in
which facility locations are integer decisions and market clearing yields
spatially resolved inherent values from LP duals. Second, we integrate three
carbon policy architectures into this model: (A) a baseline-and-credit mechanism
whose credits scale with throughput; (B) a cap-and-trade scheme in which the
carbon price is the shadow dual of the emission cap, tracing the
system's marginal abatement cost curve; and (C) a tiered carbon tax combined with
Verra VM0044-inspired removal crediting with a molar H/C $\leq$ 0.7
eligibility gate (additionality, leakage, and MRV requirements are outside the
model scope) and
grandfathering versus benchmarking allowance allocation. Third, using DOE
Billion-Ton 2023 (including its USFS FIA-based forestry module) and USDA NASS
2022 Census data, we quantify how each
policy reshapes facility location, technology selection, biochar production
(2.25--2.85 Mt/yr across the policies studied in the near-term scenario),
carbon removal (1.60--1.96 Mt
CO$_2$e/yr), and system surplus (\$0.7--3.4 billion/yr), and we show that net
emission caps looser than the system's baseline flux do not bind while removal
crediting
activates the low-price agricultural segment of demand at credit prices as low as
\$25/tCO$_2$e under the calibrated demand.

% contribution typology table (review R2-M1): literature x capability x this work
\begin{table}[htbp]
\centering
\footnotesize
\setlength{\tabcolsep}{4pt}
\caption{Model capabilities of the closest prior work versus this study:
(1)~spatially explicit network, (2)~multi-product, (3)~endogenous
price-responsive demand clearing, (4)~carbon policy embedded in the model,
(5)~carbon price obtained as a dual, (6)~multiple policy architectures compared.
A checkmark marks a capability the cited study explicitly demonstrates; a dash
marks one outside its stated scope (classification follows each work's own
framing, not an exhaustive re-analysis). No prior study embeds a carbon policy
whose price is an endogenous dual of the policy constraint.}
\label{tab:novelty}
\begin{tabular}{lcccccc}
\toprule
Study & (1) & (2) & (3) & (4) & (5) & (6)\\
\midrule
Tan ~\cite{tan2016biochar}          & \checkmark & --         & --         & --         & --         & --\\
Belmonte et al.~\cite{belmonte2018} & \checkmark & --         & --         & --         & --         & --\\
Negri et al.~\cite{negri2021}       & \checkmark & \checkmark & --         & --         & --         & --\\
Sampat et al.~\cite{sampat2019}     & \checkmark & \checkmark & \checkmark & --         & --         & --\\
Tominac et al.~\cite{tominac2022}   & \checkmark & \checkmark & \checkmark & --         & --         & --\\
\midrule
\textbf{This study} & \checkmark & \checkmark & \checkmark & \checkmark & \checkmark & \checkmark\\
\bottomrule
\end{tabular}
\end{table}

This framing separates three kinds of contribution. The
\emph{integration} contribution is the embedding of carbon policy modules
inside a spatially explicit, price-responsive, multi-product network, so that
the policy price is a model output (a dual) rather than an exogenous
parameter. The \emph{case findings} are Wisconsin-specific: the demand-bound
equilibrium at 2.30~Mt, the two-regime crossover that decides whether
crediting upgrades the pyrolysis temperature, and the two-stage response of a
net-emission cap tightened below the system's own flux. The \emph{theoretical}
contribution is the quantified boundary conditions under which a linear
emission tax is dispatch-neutral (Proposition~\ref{prop:taxinv}, evaluated
with route-specific marginal intensities in SI Section~A.11).

The paper is structured as follows. In Section~\ref{sec:model} we present the
mathematical formulation, including the base network model and the three policy
modules, together with three propositions that are proved in the
Supplementary Information: the dual-to-MAC identity of the cap (Prop.~1), the
demand-bound/supply-bound crossover prices that decide whether crediting
upgrades the pyrolysis temperature (Prop.~2), and the dispatch-neutrality
threshold of a linear emission tax in the demand-constrained regime
(Prop.~3). Section~\ref{sec:results} reports the results: baseline facility
locations and economics, environmental performance, the market-clearing
inherent values, the policy responses, the quality-differentiated extension,
and one-at-a-time sensitivities. Section~\ref{sec:discussion} discusses policy
implications against real carbon market data, and
Section~\ref{sec:conclusion} concludes; the Wisconsin case-study data and
scenario design are documented in the Supplementary Information.
